\documentclass[
    language=english,
    doctype=technical-report,
    institution=none,
    titlestyle=book,
    oneside,
    loadGlossaries,
]{../../omnilatex}

\RequirePackage{config/document-settings}
\addbibresource{../../bib/bibliography.bib}
\RequirePackage{algorithm}
\RequirePackage{algorithmic}
\RequirePackage{marginnote}
\RequirePackage{longtable}
\RequirePackage[symbol]{footmisc}
\RequirePackage[table]{xcolor}
% ─────────────────────────────────────────────────────
% Abbreviations
% ─────────────────────────────────────────────────────
\newglossaryentry{tpsatransaction}{
    type=abbreviations,
    description={Transactions Per Second --- the number of database operations completed per second},
    name={TPS},
}
\newglossaryentry{slaservicelevel}{
    type=abbreviations,
    description={Service Level Agreement --- a formal commitment between provider and client},
    name={SLA},
}
\newglossaryentry{iopsio}{
    type=abbreviations,
    description={Input/Output Operations Per Second --- a measure of storage performance},
    name={IOPS},
}
\newglossaryentry{latency}{
    type=abbreviations,
    description={The time delay between a request and its response, typically measured in milliseconds},
    name={Latency},
}
\newglossaryentry{throughput}{
    type=abbreviations,
    description={The rate at which data is transferred across a system},
    name={Throughput},
}
\newglossaryentry{nvme}{
    type=abbreviations,
    description={Non-Volatile Memory Express --- a storage protocol designed for solid-state drives},
    name={NVMe},
}
\newglossaryentry{rps}{
    type=abbreviations,
    description={Reads Per Second --- a measure of read throughput on a storage system},
    name={RPS},
}
\newglossaryentry{qla}{
    type=abbreviations,
    description={Queue Length Average --- the average number of outstanding I/O requests},
    name={QLA},
}

% ─────────────────────────────────────────────────────
% Technical Terms
% ─────────────────────────────────────────────────────
\newglossaryentry{consensus}{
    type=main,
    description={A distributed systems protocol for achieving agreement among multiple nodes},
    name={Consensus Protocol},
}
\newglossaryentry{sharding}{
    type=main,
    description={Partitioning data across multiple database instances to distribute load},
    name={Sharding},
}
\newglossaryentry{replication}{
    type=main,
    description={Copying data from one database server to another for redundancy and read scaling},
    name={Replication},
}
\newglossaryentry{connectionpool}{
    type=main,
    description={A cache of database connections maintained for reuse by applications},
    name={Connection Pool},
}

\endinput

% ─────────────────────────────────────────────────────
% Custom Commands
% ─────────────────────────────────────────────────────
\newcommand{\R}{\mathbb{R}}
\newcommand{\N}{\mathbb{N}}
\newcommand{\mat}[1]{\mathbf{#1}}
\DeclareMathOperator*{\argmin}{arg\,min}
\DeclareMathOperator*{\argmax}{arg\,max}
\DeclareMathOperator{\Tr}{Tr}
\DeclareMathOperator{\diag}{diag}

% ─────────────────────────────────────────────────────
% tcolorbox Styles
% ─────────────────────────────────────────────────────
\newtcolorbox{definitionbox}[1][]{
    colback=blue!5!white,
    colframe=blue!75!black,
    fonttitle=\bfseries,
    title=Definition,
    #1
}

\newtcolorbox{warningbox}[1][]{
    colback=red!5!white,
    colframe=red!75!black,
    fonttitle=\bfseries,
    title=Warning,
    #1
}

\newtcolorbox{notebox}[1][]{
    colback=green!5!white,
    colframe=green!50!black,
    fonttitle=\bfseries,
    title=Note,
    #1
}

\newtcolorbox{keyfindingbox}[1][]{
    colback=orange!5!white,
    colframe=orange!75!black,
    fonttitle=\bfseries,
    title=Key Finding,
    #1
}

\endinput

\renewcommand{\thefootnote}{\arabic{footnote}}

\begin{document}

\maketitle
\tableofcontents
\listoffigures
\listoftables
\listofalgorithms

\chapter{Executive Summary}

This report presents the findings of a comprehensive performance audit
conducted on Nexus Technologies' next-generation data platform during
Q1 2026. The audit scope encompassed the primary transactional database
cluster, the event-processing pipeline, the caching tier, and the
supporting network fabric across three data centres.

\begin{keyfindingbox}
The platform is currently operating at 38\% above optimal resource
utilisation. Implementation of the recommended optimisations is projected
to reduce annual infrastructure costs by approximately \$420\,000 while
improving P99 latency by 65\%.
\end{keyfindingbox}

Key findings include:
\begin{enumerate}
    \item \textbf{Storage bottleneck} --- The primary \gls{nvme} tier
        saturates at 18\,000~\gls{iopsio} during peak hours, while demand
        reaches 24\,000~IOPS (\cref{fig:iops-demand}).
    \item \textbf{Over-provisioning} --- 40\% of compute instances
        operate below 30\% utilisation (\cref{tab:compute-overview}).
    \item \textbf{Connection overhead} --- Per-query latency is elevated
        by 15\% due to suboptimal \gls{connectionpool} sizing.
    \item \textbf{Replication lag} --- Read replicas experience up to
        2.3 seconds of replication delay during peak load.
\end{enumerate}

The analysis draws on two weeks of continuous monitoring data comprising
over 12~million data points across 47~system
metrics.\footnote{All monitoring data was collected using the
open-source observability stack described in \cref{sec:monitoring}.}

\chapter{Background}

\section{System Overview}

The data platform was originally deployed in 2019 to handle an
anticipated \gls{throughput} of 10\,000~\gls{tpsatransaction}. Since
then, organic growth and three major product launches have increased peak
load to 27\,000~TPS, exceeding the original design envelope by nearly a
factor of three.\marginnote{The original capacity plan assumed 20\%
year-over-year growth. Actual growth averaged 35\% annually.}

The original deployment consisted of a three-node PostgreSQL cluster with
SAS-backed storage. In 2021, the storage tier was upgraded to SSD, and in
2023, a fourth read replica was added. However, the fundamental I/O
subsystem remained the limiting factor, as detailed
in~\cref{tab:storage-comparison}.

\section{Monitoring Methodology}
\label{sec:monitoring}

Data collection employed a distributed agent architecture with 47~metric
endpoints reporting at 60-second intervals. The monitoring stack
comprises Prometheus for time-series storage, Grafana for visualisation,
and custom Python collectors for application-level
metrics.\footnote{The collector source code is available in the
accompanying technical appendices.} All timestamps are synchronised via
NTP with sub-millisecond accuracy.

\chapter{Analysis}

\section{Performance Metrics}

\cref{tab:metrics} summarises the key performance metrics observed during
the two-week monitoring period.

\begin{table}[htbp]
    \centering
    \caption{Infrastructure performance summary.\footnotemark}
    \label{tab:metrics}
    \begin{tabular}{@{}lrrr@{}}
        \toprule
        Metric & Target & Observed & Status \\
        \midrule
        Avg.\ response time (ms) & $< 200$ & 312 & \textcolor{red}{Over} \\
        P99 \gls{latency} (ms) & $< 500$ & 1\,240 & \textcolor{red}{Over} \\
        CPU utilisation (\%) & 60--80 & 43 & \textcolor{orange}{Under} \\
        Memory utilisation (\%) & 70--85 & 62 & \textcolor{orange}{Under} \\
        Disk I/O wait (\%) & $< 10$ & 28 & \textcolor{red}{Over} \\
        \bottomrule
    \end{tabular}
\end{table}
\addtocounter{footnote}{-1}
\footnotetext{All thresholds derived from the \gls{slaservicelevel}
document dated 2024-03-15.}

The elevated disk I/O wait is the primary driver of latency degradation.
The storage subsystem saturates at approximately 18\,000~\gls{iopsio},
whereas peak demand reaches 24\,000~IOPS. \cref{fig:iops-demand}
illustrates the IOPS demand profile over a representative 24-hour period.

\begin{figure}[htbp]
    \centering
    \begin{tikzpicture}
    \begin{axis}[
        width=0.7\textwidth, height=5cm,
        xlabel={Hour}, ylabel={IOPS ($\times 10^{3}$)},
        xmin=0, xmax=24, ymin=0, ymax=30,
        grid=major, thick, legend pos=north west,
    ]
        \addplot[blue, thick, fill=blue!15] coordinates {
            (0,8) (2,5) (4,4) (6,6) (8,15) (9,22) (10,24)
            (11,23) (12,18) (13,20) (14,22) (15,21)
            (16,17) (17,14) (18,16) (19,19) (20,18)
            (21,15) (22,12) (23,10) (24,8)
        } \closedcycle;
        \addlegendentry{Demand}
        \addplot[red, dashed, thick, domain=0:24] {18};
        \addlegendentry{Capacity}
    \end{axis}
    \end{tikzpicture}
    \caption{Disk IOPS demand versus capacity over a typical day.}
    \label{fig:iops-demand}
\end{figure}

\cref{tab:storage-comparison} compares the three storage generations.

\begin{table}[htbp]
    \centering
    \caption{Storage subsystem comparison across platform generations.}
    \label{tab:storage-comparison}
    \begin{tabular}{@{}llccccc@{}}
        \toprule
        \multirow{2}{*}{Generation} & \multirow{2}{*}{Media}
            & \multicolumn{2}{c}{Sequential}
            & \multicolumn{2}{c}{Random}
            & \multirow{2}{*}{Cost/TB} \\
        \cmidrule(lr){3-4}\cmidrule(lr){5-6}
         & & Read & Write & Read & Write & \\
        \midrule
        Gen 1 (2019) & SAS HDD  & 200 MB/s & 180 MB/s & 0.2\,k & 0.15\,k & \$80 \\
        Gen 2 (2021) & SATA SSD & 550 MB/s & 520 MB/s & 50\,k   & 45\,k   & \$120 \\
        Gen 3 (2023) & \gls{nvme} SSD & 3.5 GB/s & 3.0 GB/s & 500\,k  & 400\,k  & \$200 \\
        \bottomrule
    \end{tabular}
\end{table}

\section{System Architecture}

\cref{fig:architecture} illustrates the current system architecture.

\begin{figure}[htbp]
    \centering
    \begin{tikzpicture}[
        node distance=1.2cm and 1.5cm,
        box/.style={rectangle, draw, minimum width=2.2cm, minimum height=0.8cm,
                    align=center, font=\small\sffamily},
        dataflow/.style={->, thick, >=stealth},
        highlight/.style={box, draw=red, very thick, fill=red!10},
    ]
        \node[box] (client) {Client\\Applications};
        \node[box, right=of client] (gateway) {API\\Gateway};
        \node[box, right=of gateway] (auth) {Auth\\Service};
        \node[box, right=of auth] (db) {PostgreSQL\\Cluster};
        \node[highlight, below=of db] (storage) {\gls{nvme}\\Storage};
        \node[box, below=of gateway] (cache) {Redis\\Cache};
        \node[box, below=of auth] (monitor) {Prometheus\\Monitoring};
        \draw[dataflow] (client) -- (gateway);
        \draw[dataflow] (gateway) -- (auth);
        \draw[dataflow] (auth) -- (db);
        \draw[dataflow, red, very thick] (db) -- (storage);
        \draw[dataflow] (gateway) -- (cache);
        \draw[dataflow, dashed] (monitor) -- (db);
        \draw[dataflow, dashed] (monitor) -- (cache);
    \end{tikzpicture}
    \caption{System architecture. The storage tier is the bottleneck.}
    \label{fig:architecture}
\end{figure}

\cref{fig:dataflow} shows the transformation pipeline.

\begin{figure}[htbp]
    \centering
    \begin{tikzpicture}[>=stealth', node distance=2.0cm, auto,
        proc/.style={rectangle, draw, fill=orange!15, text width=2.0cm,
                     text centered, minimum height=0.7cm, rounded corners},
    ]
        \node[proc] (ingest) {Event\\Ingestion};
        \node[proc, right=of ingest] (validate) {Schema\\Validation};
        \node[proc, right=of validate] (transform) {Data\\Transform};
        \node[proc, right=of transform] (persist) {Database\\Persist};
        \node[proc, below=of transform] (notify) {Change\\Notification};
        \draw[->, thick] (ingest) -- node {raw events} (validate);
        \draw[->, thick] (validate) -- node {validated} (transform);
        \draw[->, thick] (transform) -- node {rows} (persist);
        \draw[->, thick, dashed] (transform) -- node {CDC} (notify);
    \end{tikzpicture}
    \caption{Data flow from event ingestion to database persistence.}
    \label{fig:dataflow}
\end{figure}

\cref{fig:sequence} depicts the request lifecycle.

\begin{figure}[htbp]
    \centering
    \begin{tikzpicture}[>=stealth', every node/.style={font=\small},
        entity/.style={rectangle, draw, fill=blue!10, minimum width=1.6cm,
                       minimum height=0.6cm, align=center},
        msg/.style={->, thick},
    ]
        \node[entity] (client) at (0,0) {Client};
        \node[entity] (gw) at (3,0) {Gateway};
        \node[entity] (auth) at (6,0) {Auth};
        \node[entity] (db) at (9,0) {Database};
        \draw[dashed, gray] (client) -- (0,-5);
        \draw[dashed, gray] (gw) -- (3,-5);
        \draw[dashed, gray] (auth) -- (6,-5);
        \draw[dashed, gray] (db) -- (9,-5);
        \draw[msg] (client.north) ++(0,-0.5) -- node[above] {HTTP Request}
            (gw.north) ++(0,-0.5);
        \draw[msg] (gw.north) ++(0,-1.0) -- node[above] {Validate JWT}
            (auth.north) ++(0,-1.0);
        \draw[msg, dashed] (auth.north) ++(0,-1.3) -- node[above] {200 OK}
            (gw.north) ++(0,-1.3);
        \draw[msg] (gw.north) ++(0,-1.8) -- node[above] {Query}
            (db.north) ++(0,-1.8);
        \draw[msg, dashed] (db.north) ++(0,-2.2) -- node[above] {Result Set}
            (gw.north) ++(0,-2.2);
        \draw[msg, dashed] (gw.north) ++(0,-2.8) -- node[above] {JSON Response}
            (client.north) ++(0,-2.8);
    \end{tikzpicture}
    \caption{Request lifecycle sequence diagram.}
    \label{fig:sequence}
\end{figure}

\section{Performance Confidence Intervals}

\cref{fig:confidence} presents latency measurements with 95\% confidence
intervals.

\begin{figure}[htbp]
    \centering
    \begin{tikzpicture}
    \begin{axis}[
        width=0.8\textwidth, height=6cm,
        xlabel={Load (TPS $\times 10^{3}$)},
        ylabel={P99 \gls{latency} (ms)},
        xmin=5, xmax=30, ymin=0, ymax=1500,
        grid=major, thick, legend pos=north west,
    ]
        \addplot[blue, thick, mark=*,
            error bars/.cd, y dir=both, y explicit] coordinates {
            (5,120) +- (0,15) (10,180) +- (0,22) (15,280) +- (0,35)
            (20,520) +- (0,60) (25,940) +- (0,110) (27,1240) +- (0,150)
        };
        \addlegendentry{Measured P99}
        \addplot[red, dashed, thick, domain=5:30] {500};
        \addlegendentry{\gls{slaservicelevel} Threshold}
        \addplot[green!60!black, dotted, thick, domain=5:30] {200};
        \addlegendentry{Target P99}
    \end{axis}
    \end{tikzpicture}
    \caption{P99 latency with 95\% confidence intervals.}
    \label{fig:confidence}
\end{figure}

\section{Extended Monitoring Data}

\Cref{tab:extended-metrics} presents the full metrics broken down by
time-of-day periods.

\begin{longtable}{@{}l r r r r@{}}
    \caption{Extended performance metrics by time period.}
    \label{tab:extended-metrics} \\
    \toprule
    \textbf{Metric} & \textbf{Off-Peak} & \textbf{Business} & \textbf{Peak} & \textbf{Overall} \\
    \midrule
    \endfirsthead
    \toprule
    \textbf{Metric} & \textbf{Off-Peak} & \textbf{Business} & \textbf{Peak} & \textbf{Overall} \\
    \midrule
    \endhead
    \midrule \multicolumn{5}{r}{\textit{Continued on next page\ldots}} \\ \endfoot
    \bottomrule \endlastfoot
    Avg.\ \gls{latency} (ms)     & 145  & 298  & 487  & 312  \\
    P95 latency (ms)              & 220  & 480  & 890  & 530  \\
    P99 latency (ms)              & 380  & 920  & 1540 & 1240 \\
    \gls{throughput} (TPS)        & 4\,200 & 15\,600 & 24\,800 & 14\,900 \\
    Active connections            & 120  & 340  & 580  & 347  \\
    CPU utilisation (\%)          & 22   & 48   & 71   & 43   \\
    Memory utilisation (\%)       & 38   & 64   & 82   & 62   \\
    Disk IOPS                     & 4\,200 & 15\,800 & 24\,200 & 14\,700 \\
    Disk I/O wait (\%)            & 8    & 24   & 42   & 28   \\
    Network \gls{throughput} (Gbps) & 1.2 & 5.8  & 9.1  & 5.4  \\
    Cache hit ratio (\%)          & 94   & 87   & 72   & 84   \\
    Error rate (\textperthousand) & 0.1  & 0.3  & 1.2  & 0.5  \\
\end{longtable}

\section{Storage Health Indicators}

\Cref{tab:storage-status} provides colour-coded health indicators.

\begin{table}[htbp]
    \centering
    \caption{Storage subsystem health indicators.\footnotemark}
    \label{tab:storage-status}
    \begin{tabular}{@{}llcccc@{}}
        \toprule
        \textbf{Component} & \textbf{Metric}
            & \textbf{Normal} & \textbf{Warning} & \textbf{Critical}
            & \textbf{Current} \\
        \midrule
        \multirow{3}{*}{Primary SSD}
            & IOPS util.
            & \cellcolor{green!20}$<$60\%
            & \cellcolor{yellow!30}60--85\%
            & \cellcolor{red!20}$>$85\%
            & \cellcolor{red!30}\textbf{92\%} \\
            & Latency
            & \cellcolor{green!20}$<$5 ms
            & \cellcolor{yellow!30}5--15 ms
            & \cellcolor{red!20}$>$15 ms
            & \cellcolor{red!30}\textbf{28 ms} \\
            & Capacity
            & \cellcolor{green!20}$<$70\%
            & \cellcolor{yellow!30}70--90\%
            & \cellcolor{red!20}$>$90\%
            & \cellcolor{yellow!30}78\% \\
        \midrule
        \multirow{3}{*}{Replica SSD}
            & IOPS util.
            & \cellcolor{green!20}$<$60\%
            & \cellcolor{yellow!30}60--85\%
            & \cellcolor{red!20}$>$85\%
            & \cellcolor{green!20}34\% \\
            & Latency
            & \cellcolor{green!20}$<$5 ms
            & \cellcolor{yellow!30}5--15 ms
            & \cellcolor{red!20}$>$15 ms
            & \cellcolor{green!20}3 ms \\
            & Capacity
            & \cellcolor{green!20}$<$70\%
            & \cellcolor{yellow!30}70--90\%
            & \cellcolor{red!20}$>$90\%
            & \cellcolor{green!20}45\% \\
        \bottomrule
    \end{tabular}
\end{table}
\addtocounter{footnote}{-1}
\footnotetext{Health thresholds from vendor specs. Red cells indicate
immediate action required.}

\section{Compute Resource Overview}

\Cref{tab:compute-overview} shows compute utilisation across all instances.

\begin{table}[htbp]
    \centering
    \caption{Compute instance utilisation heatmap.\footnotemark}
    \label{tab:compute-overview}
    \begin{tabular}{@{}lccccc@{}}
        \toprule
        \textbf{Instance} & \textbf{CPU (\%)} & \textbf{Mem (\%)}
            & \textbf{IOPS} & \textbf{Network} & \textbf{Status} \\
        \midrule
        db-primary-01
            & \cellcolor{red!30}92 & \cellcolor{orange!30}78
            & \cellcolor{red!30}24\,200 & \cellcolor{yellow!30}8.1
            & \textcolor{red}{Critical} \\
        db-primary-02
            & \cellcolor{orange!30}71 & \cellcolor{yellow!30}64
            & \cellcolor{orange!30}15\,800 & \cellcolor{green!20}4.2
            & \textcolor{orange}{Warning} \\
        db-replica-01
            & \cellcolor{green!20}34 & \cellcolor{green!20}45
            & \cellcolor{green!20}8\,200 & \cellcolor{green!20}2.1
            & Healthy \\
        db-replica-02
            & \cellcolor{green!20}28 & \cellcolor{green!20}38
            & \cellcolor{green!20}6\,400 & \cellcolor{green!20}1.8
            & Healthy \\
        cache-node-01
            & \cellcolor{yellow!30}55 & \cellcolor{green!20}42
            & --- & \cellcolor{green!20}3.5 & Healthy \\
        cache-node-02
            & \cellcolor{green!20}48 & \cellcolor{green!20}40
            & --- & \cellcolor{green!20}3.2 & Healthy \\
        \bottomrule
    \end{tabular}
\end{table}
\addtocounter{footnote}{-1}
\footnotetext{Colour thresholds: green $<$60\%, yellow 60--80\%, orange
80--90\%, red $>$90\%.}

\section{Resource Allocation Algorithm}

\Cref{alg:optimise} describes the greedy algorithm for optimal resource
allocation.

\begin{algorithm}[htbp]
\caption{Greedy resource allocation optimiser.}
\label{alg:optimise}
\begin{algorithmic}[1]
\REQUIRE Workload set $W = \{w_1, \ldots, w_n\}$, resource pool $R$
\ENSURE  Optimal allocation $A^*$
\STATE $A \gets$ \textsc{InitialAllocation}$(W, R)$
\STATE $\Delta \gets \infty$
\WHILE{$\Delta > \epsilon$}
    \STATE $\Delta \gets 0$
    \FORALL{$w_i \in W$}
        \STATE $r^* \gets \argmin_{r \in R} \textsc{Cost}(w_i, r)$
        \IF{$\textsc{Cost}(w_i, r^*) < \textsc{Cost}(w_i, A(w_i))$}
            \STATE $A(w_i) \gets r^*$
            \STATE $\Delta \gets \Delta + 1$
        \ENDIF
    \ENDFOR
\ENDWHILE
\RETURN $A$
\end{algorithmic}
\end{algorithm}

\begin{definitionbox}[title=Complexity Analysis]
The algorithm converges in $O(n \cdot |R|)$ per iteration, with
empirical convergence in fewer than 5 iterations. Total complexity is
$O(k \cdot n \cdot |R|)$, where $k$ is the number of iterations.
\end{definitionbox}

\section{Performance Comparison}

\Cref{fig:perf-comparison} compares throughput before and after
optimisations.

\begin{figure}[htbp]
    \centering
    \begin{tikzpicture}
    \begin{axis}[
        width=0.85\textwidth, height=6cm,
        ybar, bar width=14pt,
        ylabel={Throughput (TPS)},
        symbolic x coords={Ingestion, Transform, Persist, Query},
        xtick=data, ymin=0, ymax=35000,
        legend style={at={(0.5,-0.18)}, anchor=north, legend columns=2, font=\small},
        grid=major, grid style={dashed, gray!30},
        nodes near coords,
        every node near coord/.append style={font=\tiny, rotate=90, anchor=west},
    ]
        \addplot[fill=blue!60] coordinates
            {(Ingestion,27000) (Transform,22000) (Persist,18000) (Query,15000)};
        \addplot[fill=green!60] coordinates
            {(Ingestion,30000) (Transform,28000) (Persist,27000) (Query,25000)};
        \addlegendentry{Current}
        \addlegendentry{Projected}
    \end{axis}
    \end{tikzpicture}
    \caption{Throughput: current vs.\ projected after optimisations.}
    \label{fig:perf-comparison}
\end{figure}

\section{Resource Utilisation Over Time}

\Cref{fig:resource-util} shows stacked utilisation over 24 hours.

\begin{figure}[htbp]
    \centering
    \begin{tikzpicture}
    \begin{axis}[
        width=0.85\textwidth, height=6cm,
        xlabel={Hour}, ylabel={Utilisation (\%)},
        xmin=0, xmax=24, ymin=0, ymax=100,
        grid=major, thick, legend pos=north west, area style,
    ]
        \addplot[fill=blue!40, draw=blue!60, thick, stack plots=y]
            coordinates {(0,22) (4,15) (6,18) (8,35) (10,48)
                         (12,42) (14,48) (16,38) (18,32) (20,28)
                         (22,24) (24,22)};
        \addlegendentry{CPU}
        \addplot[fill=green!40, draw=green!60, thick, stack plots=y]
            coordinates {(0,38) (4,32) (6,35) (8,55) (10,72)
                         (12,68) (14,72) (16,60) (18,52) (20,45)
                         (22,40) (24,38)};
        \addlegendentry{Memory}
        \addplot[fill=orange!40, draw=orange!60, thick, stack plots=y]
            coordinates {(0,8) (4,4) (6,12) (8,32) (10,42)
                         (12,36) (14,40) (16,30) (18,22) (20,16)
                         (22,12) (24,8)};
        \addlegendentry{I/O Wait}
    \end{axis}
    \end{tikzpicture}
    \caption{Stacked resource utilisation over 24 hours.}
    \label{fig:resource-util}
\end{figure}

\section{Latency Distribution}

\Cref{fig:latency-hist} presents the histogram of response latencies.

\begin{figure}[htbp]
    \centering
    \begin{tikzpicture}
    \begin{axis}[
        width=0.8\textwidth, height=6cm,
        xlabel={Latency (ms)}, ylabel={Frequency},
        ymin=0, ybar, bar width=8pt,
        grid=major, grid style={dashed, gray!30},
    ]
        \addplot[fill=blue!50, draw=blue!70] coordinates {
            (50,120) (100,340) (150,580) (200,720) (250,540)
            (300,380) (350,220) (400,150) (450,90) (500,60)
            (600,35) (800,18) (1000,8) (1200,3)
        };
        \addplot[red, thick, domain=0:1300, samples=100]
            {800*exp(-(x-200)^2/20000)};
    \end{axis}
    \end{tikzpicture}
    \caption{Response latency histogram with fitted normal curve.}
    \label{fig:latency-hist}
\end{figure}

\section{Configuration and Code}

\Cref{code:config} shows the monitoring agent configuration.

\begin{listing}[htbp]
    \centering
    \inputminted[fontsize=\scriptsize, linenos, frame=lines, framesep=2mm]
        {ini}{config/monitoring.ini}
    \caption{Monitoring agent configuration file.}
    \label{code:config}
\end{listing}

\Cref{code:api} presents the REST API endpoint definitions.

\begin{listing}[htbp]
    \centering
    \inputminted[fontsize=\scriptsize, linenos, frame=lines, framesep=2mm]
        {http}{config/api-routes.http}
    \caption{Metrics API endpoint definitions.}
    \label{code:api}
\end{listing}

\Cref{code:sql} shows the SQL queries for aggregate metric computation.

\begin{listing}[htbp]
    \centering
    \inputminted[fontsize=\scriptsize, linenos, frame=lines, framesep=2mm]
        {sql}{config/queries.sql}
    \caption{Aggregate metric computation queries.}
    \label{code:sql}
\end{listing}

The agents use \mintinline{python}{asyncio} for concurrent data
collection, with configuration from \mintinline{ini}{monitoring.ini}.
Results are stored in \mintinline{sql}{PostgreSQL} with
\mintinline{python}{pool_size = 20}.\marginnote{Connection pool sizing
should be validated under peak load.}

\chapter{Recommendations}
\label{sec:recommendations}

\begin{notebox}
Recommendations ordered by impact. Full details in \cref{app:requirements}.
\end{notebox}

Based on \cref{chap:analysis}, we recommend:
\begin{enumerate}
    \item \textbf{Migrate to \gls{nvme}-backed storage} --- Raises the
        IOPS ceiling to 80\,000. Timeline: 4 weeks.
    \item \textbf{Right-size compute instances} --- Reducing CPU and
        memory by 30--35\% on 40\% of workloads. Annual savings: \$96\,000.
    \item \textbf{Implement \gls{connectionpool} tuning} --- Reduces
        per-query overhead by an estimated 15\%.
    \item \textbf{Introduce read replicas} --- Offloads reporting queries
        to dedicated replicas.
\end{enumerate}

\printglossary[type=abbreviations,title={List of Abbreviations}]
\printglossary[type=main,title={List of Technical Terms}]

\printbibliography

\appendix

\chapter{System Requirements}
\label{app:requirements}

\begin{table}[htbp]
    \centering
    \caption{Minimum hardware requirements for recommended deployment.}
    \label{tab:hardware-reqs}
    \begin{tabular}{@{}lccc@{}}
        \toprule
        Component & Specification & Quantity & Notes \\
        \midrule
        CPU        & 16-core, 3.2 GHz   & 4 nodes & AMD EPYC or Intel Xeon \\
        Memory     & 128 GB DDR4 ECC     & 4 nodes & Minimum per node \\
        Storage    & 2 TB \gls{nvme} Gen4      & 4 nodes & RAID-10 \\
        Network    & 25 GbE              & 8 ports & Dual-homed per node \\
        \bottomrule
    \end{tabular}
\end{table}

Software: Ubuntu 22.04 LTS or RHEL 9; PostgreSQL 15+; Redis 7.0+;
Prometheus 2.45+, Grafana 10.0+; Docker 24.0+.

\chapter{API Reference}
\label{app:api}

Ingest metrics:
\begin{minted}[fontsize=\small]{http}
POST /api/v2/metrics/ingest
Content-Type: application/json
Authorization: Bearer <token>

{"hostname":"db-primary","timestamp":"2026-03-15T14:30:00Z",
 "metrics":{"latency_ms":312,"iops":24000,"cpu_percent":71,"memory_percent":82}}
\end{minted}

Query metrics:
\begin{minted}[fontsize=\small]{http}
GET /api/v2/metrics/query?host=db-primary&window=2h&percentile=99
Authorization: Bearer <token>
\end{minted}

Response:
\begin{minted}[fontsize=\small]{json}
{"host":"db-primary","window":"2h","percentile":99,
 "value":1240.0,"unit":"ms","sample_count":240000}
\end{minted}

\chapter{Test Results}
\label{app:test-results}

\begin{table}[htbp]
    \centering
    \caption{Load test results under recommended configuration.}
    \label{tab:load-tests}
    \begin{tabular}{@{}lrrrrr@{}}
        \toprule
        Test & TPS & Avg.\ (ms) & P95 (ms) & P99 (ms) & Errors \\
        \midrule
        Baseline        & 10\,000 & 85  & 142 & 210 & 0 \\
        Normal load     & 15\,000 & 120 & 198 & 310 & 0 \\
        Peak load       & 25\,000 & 180 & 320 & 480 & 2 \\
        Stress test     & 30\,000 & 290 & 580 & 890 & 18 \\
        Endurance (4h)  & 15\,000 & 118 & 195 & 298 & 0 \\
        \bottomrule
    \end{tabular}
\end{table}

The P99 latency distribution was analysed using Kolmogorov--Smirnov
testing: $D_{240000} = 0.034$ with $p < 0.001$, indicating departure
from log-normality.

\chapter{Glossary Reference}
\label{app:glossary}

\begin{table}[htbp]
    \centering
    \caption{Technical terms quick reference.}
    \label{tab:glossary-summary}
    \begin{tabular}{@{}lp{0.6\textwidth}@{}}
        \toprule
        \textbf{Term} & \textbf{Definition} \\
        \midrule
        \gls{tpsatransaction} & Database operations per second \\
        \gls{iopsio} & Storage I/O operations per second \\
        \gls{latency} & Request-response time delay \\
        \gls{throughput} & Data transfer rate \\
        \gls{nvme} & Storage protocol for solid-state drives \\
        \gls{slaservicelevel} & Provider-client commitment \\
        \gls{connectionpool} & Cached database connections \\
        \gls{consensus} & Distributed agreement protocol \\
        \gls{sharding} & Data partitioning across instances \\
        \gls{replication} & Data copying for redundancy \\
        \bottomrule
    \end{tabular}
\end{table}

\end{document}
